\documentclass[10pt,a4paper]{amsart}
\usepackage[margin=19mm]{geometry}
\usepackage{amsmath,amssymb,mathtools}
\usepackage[scr=rsfs,cal=euler]{mathalfa}
\usepackage{xfrac}
\usepackage[unicode,hidelinks,bookmarks=false]{hyperref}
\newcommand{\ie}{i.e.\ }
\newcommand{\cf}{cf.\ }
\newcommand{\resp}{respectively\ }
\newcommand{\Subsection}{\subsection}
\providecommand{\nobreakdash}{\nobreak\hbox{-}}
\setlength{\emergencystretch}{2em}
\multlinegap0pt
\usepackage{bm}
\usepackage{bbm}
\usepackage{dsfont}
\usepackage{xspace}
\usepackage{cancel}
\usepackage{mathtools}
\usepackage{amsthm}
\makeatletter
\@ifundefined{theorem}{\newtheorem{theorem}{Theorem}}{}
\@ifundefined{lemma}{\newtheorem{lemma}[theorem]{Lemma}}{}
\makeatother
\newcommand{\dx}{\partial_x}
\title[The traveling wave solutions of the 1D hyperbolic Keller-Seg]{The traveling wave solutions of the 1D hyperbolic Keller-Segel equations}
\author{Xin Liu, William Kyle Barker}
\date{Source: arXiv:2608.12226v1 (CC BY 4.0)}
\begin{document}
\maketitle

\noindent\emph{Source and licence.} The traveling wave solutions of the 1D hyperbolic Keller-Segel equations. Version arXiv:2608.12226v1 (CC BY 4.0), distributed under the Creative Commons Attribution 4.0 International licence, as recorded in the arXiv licence metadata for this exact version and shown on the abstract page https://arxiv.org/abs/2608.12226v1. Full rights notice: licenses/arxiv-2608.12226-CC-BY-4.0.txt.\par\smallskip

\noindent\textit{Editorial scope.} Counted unit: the Lemma whose source label is \texttt{lm:no-crossing-asymptotics} in the article (source0.tex lines 2434--2440, source label \texttt{lm:no-crossing-asymptotics}), delivered together with the article's own complete original proof of it (source0.tex lines 2442--2473, one \texttt{proof} environment of 32 lines). No mathematics is written, repaired or re-derived: statement and proof are the article's own text line for line. Required context retained uncounted: eq:jump-1 (lines 243--250); sys:tw (lines 702--714); eq:tw-02 (lines 704--713); eq:tw-nv-ff (lines 2338--2341); eq:tw-nv-001 (lines 2343--2346); lm:no-ttrivial-point (lines 2380--2388); eq:tw-nv-003 (lines 2383--2386); eq:tw-nv-005 (lines 2400--2404). Every definition, display and label of those lines is retained unchanged. Omitted from this package and not redistributed: 1--242, 251--701, 714--2337, 2342--2342, 2347--2379, 2387--2399, 2405--2433, 2441--2441, 2474--2853. Nothing inside a retained excerpt is omitted or altered: every retained body is delivered exactly as the article's lines are written, so no omission receipt of any kind accompanies this package. Citations: the retained excerpts carry no citation command at all, so no bibliography entry of the article is transcribed and no reference is redistributed. Reference adaptation: none was needed and none was made; every reference inside the delivered unit resolves inside it, so no unresolved reference or citation is produced. Printed slips kept literal: no printed word, symbol, hypothesis or number of the counted statement or of its proof was altered. Layout and class: the article's own class is \texttt{article} with packages including amsmath, amssymb, amsthm, xcolor, graphicx, hyperref, geometry, pgfplots; the wrapper is the lane's pinned \texttt{amsart} editorial wrapper at 10pt on A4, which restates the theorem-like environments the retained text needs and adds the article's own macro definitions that the retained text needs (\texttt{\textbackslash dx}), with extra wrapper packages (bm, bbm, dsfont, xspace, cancel, mathtools). The delivered numbering is the unit's own. Assets: no retained span contains a figure environment, a graphics-inclusion command, a PDF-page inclusion, a table or a TikZ picture; no image file is copied into this package, the package declares no asset dependency and no image file is redistributed with it. Licence: version arXiv:2608.12226v1 of this article is distributed under the Creative Commons Attribution 4.0 International licence, as recorded in the arXiv licence metadata for this exact version and shown on the abstract page https://arxiv.org/abs/2608.12226v1. A full rights notice is delivered at licenses/arxiv-2608.12226-CC-BY-4.0.txt. Characters: every retained excerpt is pure ASCII, so the delivered engine, XeLaTeX with its default Latin Modern text font, reports no missing character.
    \begin{align}
        [[S]] = [[\dx S]] & = 0, \label{eq:jump-1} \\
        x'(t) & = \dx S(1- \sigma^+ - \sigma^-), \label{eq:jump-2} \qquad \text{and} \\
        \label{eq:jump-3}& \begin{cases}
            \dx S \leq 0 & \text{if} \ \sigma^+ > \sigma^-,\\
            \dx S \geq 0 & \text{if} \ \sigma^+ < \sigma^-.
        \end{cases}
    \end{align}

\begin{subequations}
    \label{sys:tw}
\begin{align}
-c\partial_\xi\sigma_c
+\partial_\xi\big(\sigma_c(1-\sigma_c)\partial_\xi S_c\big)
&=0,
\label{eq:tw-01}
\\[4pt]
-\partial_{\xi\xi}S_c+S_c
&=\sigma_c.
\label{eq:tw-02}
\end{align}
\end{subequations}

\begin{equation}
    \label{eq:tw-nv-ff}
    \sigma_c(\xi), \ S_c(\xi) \rightarrow \sigma_\infty, \qquad \text{as} \ \xi \rightarrow - \infty. 
\end{equation}

\begin{equation}
    \label{eq:tw-nv-001}
    - c \sigma_c + \sigma_c (1-\sigma_c) \partial_\xi S_c = - c \sigma_\infty.
\end{equation}

\begin{lemma}
\label{lm:no-ttrivial-point}
    If there exists $ \xi_0 \in \mathbb R $ such that $ S_c(\xi_0) = \sigma_c(\xi_0) = \sigma_\infty $, then 
    \begin{equation}
        \label{eq:tw-nv-003}
        \sigma_c = S_c = \sigma_\infty \qquad \forall \ \xi \in \mathbb R. 
    \end{equation}
    We call \eqref{eq:tw-nv-003} the trivial solution to \eqref{sys:tw} with \eqref{eq:tw-nv-ff}.
\end{lemma}

    \begin{equation}
    \label{eq:tw-nv-005}
        \partial_\xi R_c = T_c  = \frac{c \eta_c}{\sigma_c(1-\sigma_c)}, \qquad \partial_\xi T_c %= S_c - \sigma_c = R_c - \eta_c 
        = R_c - \frac{\sigma_c (1-\sigma_c)}{c}T_c, 
    \end{equation}

\begin{lemma}
    \label{lm:no-crossing-asymptotics}
    Other than the trivial solution \eqref{eq:tw-nv-003}, there is no solution $ (\sigma_c, S_c) $ to \eqref{sys:tw} with \eqref{eq:tw-nv-ff} admitting a point $ \sigma_c (\xi_0) = \sigma_\infty $ for some $ \xi_0 \in \mathbb R $.

    In other words, 
    if there exists a point $ \xi_0 \in \mathbb R $ such that $ \sigma_c(\xi_0) = \sigma_\infty $, then $ (\sigma_c,S_c) $ is the trivial solution \eqref{eq:tw-nv-003}. 
\end{lemma}

\begin{proof}
    We prove by contradiction. Suppose that the set 
    \begin{equation}
        \label{001-tw-nv}
        \mathsf M: = \lbrace \xi \vert \sigma_c(\xi) = \sigma_\infty \rbrace \neq \emptyset. 
    \end{equation}
    Let 
    \begin{equation}
        \label{002-tw-nv}
        \xi_0 = \inf \mathsf M. 
    \end{equation}
    
    First, we claim that $ \xi_0 > -\infty $. Otherwise, there exists $ \xi_1 \in \mathbb R $ such that for all $ \xi < \xi_1 $,  $ \sigma_c (\xi) = \sigma_\infty $ and $ S_c (\xi) = \sigma_\infty $ thanks to \eqref{eq:tw-nv-001} and \eqref{eq:tw-nv-ff}. Lemma \ref{lm:no-ttrivial-point} implies that $ (\sigma_c,S_c) $ is a trivial solution. 
    
    Next, without loss of generality, assume that $ \sigma_c > \sigma_\infty $ for $ \xi < \xi_0 $ and $ c > 0 $. Then from \eqref{eq:tw-nv-001} (or equivalently \eqref{eq:tw-nv-005}), one has that
    \begin{equation}
        \label{003-tw-nv}
        \partial_\xi S_c = \frac{c (\sigma_c - \sigma_\infty)}{\sigma_c (1-\sigma_c)} > 0 \quad \forall \ \xi < \xi_0 \qquad \text{and} \qquad  \partial_\xi S_c(\xi_0) = 0.
    \end{equation}
    Hence $ S_c(\xi) > S_c(-\infty) = \sigma_\infty $ for all $ \xi \leq  \xi_0 $. However, from \eqref{eq:tw-02}, one has that
    \begin{equation}
    \label{004-tw-nv}
        \partial_{\xi\xi} S_c (\xi_0) = S_c(\xi_0) - \sigma_c(\xi_0) > 0.
    \end{equation}
    Thanks to \eqref{eq:jump-1} and \eqref{eq:tw-nv-001}, $ S_c $ and $ \sigma_c $ are continuous at $ \xi_0 $. Thus for some $ \xi_\delta < \xi_0 $, one can calculate that
    \begin{equation}
         0 < \frac{c (\sigma_c - \sigma_\infty)}{\sigma_c (1-\sigma_c)}\Big\vert_{\xi = \xi_1} = \partial_\xi S_c(\xi_1) = - \int_{\xi_1}^{\xi_0} \partial_{\xi\xi} S_c(\xi') \, d\xi' < 0,
    \end{equation}
    which is impossible.
    The case when $ c > 0 $ and/or $ \sigma_c(\xi) < \sigma_\infty $ for $\xi< \xi_0 $ follows with similar arguments. 
    This finishes the proof of the lemma. 
\end{proof}

\end{document}
